\documentclass[11pt]{article}

\usepackage[review]{acl}

\usepackage{times}
\usepackage{latexsym}
\usepackage[T1]{fontenc}
\usepackage[utf8]{inputenc}
\usepackage{microtype}
\usepackage{inconsolata}
\usepackage{graphicx}

\title{Project Proposal: Grapheme-to-Phoneme Mapping for Kashmiri Language (Sharada Script)}

\begin{document}
\maketitle


%------------------------------------------------
\section{Proposal}

\textbf{Problem Statement.}
Grapheme-to-phoneme (G2P) conversion is a fundamental task in computational linguistics that maps written characters to their corresponding pronunciations. Kashmiri, particularly in the historical Sharada script, lacks computational resources and automated pronunciation tools. This significantly limits the development of speech technologies and creates a strong need for an automated G2P system for the Kashmiri language.

\textbf{Motivation.}
Most existing G2P systems focus on high-resource languages such as English or Hindi due to their strong digital presence. Kashmiri, despite being a scheduled language with cultural and linguistic importance, remains computationally underrepresented. Developing a G2P system for Sharada-script Kashmiri contributes toward language preservation while enabling speech synthesis and linguistic analysis tools for low-resource languages.

\textbf{Proposed Approach.}
We propose building a hybrid rule-based and data-driven Grapheme-to-Phoneme model. The workflow consists of the following steps:

\begin{itemize}
    \item \textbf{Text Extraction:} Collect Kashmiri text from manuscripts, archives, and linguistic resources. Since Kashmiri has limited digital availability, extraction from multiple physical and digital sources will be required.
    
    \item \textbf{Normalization:} Perform preprocessing including removal of noise and punctuation, script cleaning, and preparation of Unicode-compliant Kashmiri text.
    
    \item \textbf{Transliteration:} Convert normalized Kashmiri text into Roman and International Phonetic Alphabet (IPA) representations for compatibility with speech processing tools.
    
    \item \textbf{Speech Generation:} Reintroduce punctuation required for natural prosody and pass IPA sequences into an IPA-to-speech conversion system.
\end{itemize}


%------------------------------------------------
\section{Dataset}

\textbf{Source.}
Text data will be collected from publicly available Kashmiri linguistic resources, digitized Sharada manuscripts, and transliteration datasets. OCR techniques will also be used to digitize handwritten and printed materials.

\textbf{Description.}
The dataset will contain Kashmiri words written in Sharada script paired with phonemic representations using IPA or standardized phoneme notation.

\textbf{Relevance.}
Aligned grapheme–phoneme data directly supports pronunciation modeling. OCR-based digitization enables utilization of otherwise inaccessible resources for this low-resource language.

\textbf{Preprocessing.}
Processing steps include script normalization, Unicode standardization, grapheme segmentation, phoneme alignment, and dataset splitting.

\textbf{Availability.}
Although Kashmiri is a low-resource language, substantial non-digitized textual material exists in manuscripts and books which will be converted into usable datasets.


%------------------------------------------------
\section{Deliverables}

The expected outcomes of this project include:

\begin{itemize}
\item Kashmiri Unicode-based transliteration system to Roman and IPA
\item Curated Sharada grapheme--phoneme dataset generated via OCR
\item Rule-based baseline G2P conversion system
\item Kashmiri text-to-speech capability from non-digital sources
\item Evaluation using Phoneme Error Rate (PER)
\item Source code and documentation
\item Extendable framework for future phoneme-to-grapheme conversion
\end{itemize}


%------------------------------------------------
\section{Proposed Timeline}
Our project will be executed over a 6-week period, structured around the following concrete milestones shown in Table 1. Note that this project uses simple and straight forward technical aspencts, but emphasizes more on the handling on linguisctic edge cases, and ability to run on as many Kashmiri documents as possible.
Time have been allocated accordingly.  
\begin{table}[t]

\small
\setlength{\tabcolsep}{3pt}

\begin{tabular}{p{2.6cm} p{3.6cm} p{3.8cm}}
\hline
\textbf{Period} & \textbf{Task Completed} & \textbf{Milestone} \\
\hline

Week 1 (Feb End)
& Collection of Sharada manuscripts
& Raw Kashmiri textual resources gathered \\

Week 2 (Early March)
& OCR processing and digitization
& Machine-readable corpus created \\

Week 3 (Mid March)
& Text normalization and Unicode setup
& Clean Unicode-ready dataset \\

\textbf{23 March}
& \textbf{Interim Presentation}
& \textbf{Normalized corpus demonstrated} \\

Week 4 (Late March)
& Grapheme preprocessing
& G2P-ready structured input \\

Week 5 (Early April)
& Roman and IPA transliteration with speech integration
& Working G2P + speech pipeline \\

Week 6 (Mid April)
& Testing, evaluation and PER calculation
& Optimized pronunciation system \\

\textbf{21 April}
& \textbf{Final Submission}
& \textbf{Complete G2P speech system} \\

\hline
\end{tabular}

\caption{Project timeline showing completed tasks and achieved system milestones}
\end{table}

\end{document}
